\documentclass [a4paper,10pt] {article}

\usepackage{graphicx}
\usepackage{german}
\usepackage{titlesec}

\titleformat*{\section}{\LARGE \bfseries}
\titleformat*{\subsection}{\Large \bfseries}

\titlespacing*{\section}{0pt}{*4}{*2}
\titlespacing*{\subsection}{0pt}{*4}{*1.5}

\setlength{\parindent}{0pt}

\begin {document}

\begin{titlepage}
	\raggedleft
	\includegraphics[width=5cm]{HFULogo}\par\vspace{1cm}

	\centering
	\vfill

	\LARGE FHF Train \par
	\vspace{1cm}
	
	\Large Semesterprojekt 2018 \\
	\vfill
	
	\Large Autoren: \\ 
	Stephan Mosbach AIB 6 \\
 	Ilirian Halilaj AIB6 \\ 
	Jonas Wiedle AIN6 \\ 
	Kevin Junk AIN4 \\
	Alexander Leibinger AIN6\par
	
	\vfill
	Referent: Prof. Dr. Rainer M"uller\par

	\vfill
\end{titlepage}

	\section*{Abstract}

	\vspace{5mm}

	This Project is presented by the University of Furtwangen.

\newpage

\tableofcontents

\newpage
	\section{Einleitung}
	
		Mit dieser Dokumentation wollen wir den Lesern das Vorgehen unserer Projektgruppe für das Semesterprojekt „FHFTrain“ n"aherbringen und einen Einblick in die 			
		Problematik und die daraus resultierenden Ergebnisse im Rahmen des Projektes aus einer projektorganisatorischen und einer informationstechnischen Sicht darstellen. \par
		Das Ziel unseres Projektes ist es ein Yocto Image von einem Server/Raspberry Pi über das Netzwerk automatisch auf den Zug geladen und ausgeführt werden.
	
\newpage
	\section{Projektbeschreibung}
\newpage
	\section{Zeitplan}
		\vfill
		\subsection{gesch"atzter Zeitplan}
		\vfill
		\subsection{Zeitplan Halbzeit}
		\vfill
		\subsection{endg"ultiger Zeitplan}
		\vfill



\end {document}


